\documentclass[11pt]{../../modelos-latex/classes/ifscarticle}
\usepackage{../../modelos-latex/classes/ifscutils}
\geometry{a4paper,hmargin=2cm,top=3.5cm,bottom=2cm,headheight=3cm,heightrounded}

\usepackage[style=abnt,noslsn,justify]{biblatex}
\usepackage{csquotes}
% Bibliografia do(a) aluno(a): exporte o seu .bib do Zotero/Mendeley com o
% nome referencias.bib e deixe-o na mesma pasta deste arquivo. O arquivo
% distribuído junto com este template já traz entradas de exemplo.
\addbibresource{referencias.bib}

% -----------------------------------------------------------------
% Template de resumo estendido -- TCC1 EngTel
% Estrutura conforme wiki de recursos da UC (TCC1-EngTel, página):
%   resumo <=250 palavras; introdução/metodologia/resultados <=1000
%   palavras cada; considerações <=200 palavras; 3-5 palavras-chave.
% -----------------------------------------------------------------

\begin{document}

\begin{center}
	{\Large \textbf{Autenticação e delegação dinâmica em ecossistema agêntico
}}\vspace{.3cm}

	Marcos Vinicius Wagner\textsuperscript{1}, Emerson Ribeiro de Mello\textsuperscript{2}

	\vspace{.2cm}
	{\small
	\textsuperscript{1}Curso de Engenharia de Telecomunicações, IFSC Câmpus São José\\
	\textsuperscript{2}Orientador(a), IFSC Câmpus São José
	}
\end{center}

\vspace{.5cm}

\section*{Resumo}
% Máximo de 250 palavras. Contextualize o tema, o problema, o objetivo e a
% abordagem esperada em um único parágrafo.
O avanço da Inteligência Artificial agêntica introduziu sistemas capazes de atuar com alto grau de autonomia, interagindo com ferramentas externas e outros agentes através de diferentes domínios de confiança. No entanto, essa autonomia traz um desafio. É necessário identificar esse agente. Quando um serviço receber uma solicitação dele, o serviço deve entender que a ação está sendo realizada em nome do usuário. A ausência de uma arquitetura de identidade dedicada gera falhas graves de segurança, como a personificação. Nesse contexto, este trabalho tem como objetivo analisar e estruturar modelos arquiteturais de gestão de identidade para agentes autônomos. A pesquisa será desenvolvida por meio de uma revisão da literatura sobre agentes de IA e gestão de identidade e acesso, seguida da identificação de requisitos e da análise comparativa de abordagens arquiteturais existentes. Como resultado, espera-se elaborar uma proposta de modelo arquitetural de gestão de identidade para agentes de IA. Espera-se, assim, contribuir para a discussão e estruturação de mecanismos de identidade e delegação adequados à atuação autônoma de agentes de IA.

\noindent\textbf{Palavras-chave:} IA agêntica; gestão de identidade e acesso; Segurança da informação.
% 3 a 5 palavras-chave, separadas por ponto-e-vírgula.

\section{Introdução}
% Máximo de 1000 palavras. Contextualização do tema, motivação, problema de
% pesquisa e objetivos (geral e específicos).
O avanço da Inteligência Artificial agêntica introduziu sistemas capazes de atuar com alto grau de autonomia, interagindo com ferramentas externas e outros agentes através de diferentes domínios de confiança. É necessário identificar esse agente. A ausência de uma arquitetura de identidade dedicada gera falhas graves de segurança, como a personificação, um cenário em que o agente de IA utiliza as credenciais do usuário de forma indistinguível, como login e senha. Isso impossibilita que os logs de auditoria diferenciem as ações humanas das ações da máquina, levando à perda de rastreabilidade, e perda de limites das ações que a IA pode executar. O serviço deve entender que a ação está sendo realizada em nome do usuário, mas entendendo que está sendo executada por um agente.

É fundamental para essa resolução resgatar o conceito de Gestão de Identidade e Acesso (GID). A GID busca garantir que as entidades certas acessem os recursos certos, pelas razões certas. Esse controle é sustentado por três pilares chamados de AAA, a Autenticação (verificação de quem é a entidade), a Autorização (controle de quais privilégios a entidade possui) e a Auditoria (rastreabilidade das ações executadas pela entidade). \cite{melloGestaoIdentidadeAcesso}

Em modelos com agentes autônomos, o grande desafio é fazer com que esses três pilares operem de forma segura sem a intervenção humana constante. Isso exige uma mudança de paradigma da ``personificação'' para a ``delegação'', também conhecida como padrão \textit{On-Behalf-Of} (OBO)\cite{openid-agentic-identity}. Nesse modelo, a infraestrutura de identidade passa a distinguir claramente o usuário que delega a tarefa, do agente de IA que a executa, garantindo que os protocolos de transporte carreguem ambas as identidades de forma vinculada e assinada.

\subsection{Objetivo Geral}
O objetivo geral deste trabalho é analisar e estruturar modelos arquiteturais de gestão de identidade para IA agêntica, baseados no padrão de delegação OBO, para garantir a rastreabilidade, o controle de escopo e a segurança nas interações entre múltiplos domínios de confiança. 

\begin{itemize}
    \item \textbf{Estudar as abordagens de personificação e delegação:} Estabelecer os requisitos para que a infraestrutura possa diferenciar e vincular a identidade do usuário que delegou e a do agente que está executando.
    
    \item \textbf{Comparar os modelos arquiteturais de identidade:} Avaliar implementações de delegação para agentes baseado em seus modelos de identificadores, protocolos de delegação e infraestruturas.
\end{itemize}


\section{Metodologia}
% Máximo de 1000 palavras. Método de trabalho previsto para investigar/atacar
% o problema. Em TCC1 pode ser uma metodologia preliminar, a ser detalhada no
% documento de pré-projeto.

O trabalho será desenvolvido por meio de uma pesquisa dividida em etapas. Inicialmente, será realizada uma revisão da literatura sobre Agentes de IA para entender suas principais arquiteturas, cenários de uso e tecnologias envolvidas. Então, uma revisão focada na gestão de identidade e acesso aplicada a agentes de IA, com foco nos mecanismos de autenticação, autorização, delegação e auditoria. Serão investigadas abordagens que permitam distinguir a identidade do usuário da identidade do agente que atua em seu nome, e mecanismos de identidade e credenciais digitais para ambientes distribuídos.

Serão identificados os principais requisitos para a gestão de identidade em sistemas multiagênticos. Entre os requisitos considerados estão a identificação individual do agente, a vinculação entre o agente e o usuário que delegou a ação, a limitação do escopo das permissões, a possibilidade de rastrear as ações realizadas, a capacidade de operar entre diferentes domínios, e o ciclo de vida dos agentes.

Serão selecionadas e analisadas as possíveis abordagens arquiteturais de gestão de identidade para agentes autônomos, considerando os requisitos identificados. Os resultados serão organizados de forma comparativa, destacando os potenciais e limitações de cada abordagem. 

Será elaborado um levantamento de requisitos e um modelo comparativo de diferentes abordagens de gestão de identidade para agentes, considerando os requisitos estudados. Será estudado um caso de uso, no qual serão comparados os potenciais e limitações de cada abordagem.

\section{Resultados e discussão esperados}
% Máximo de 1000 palavras. O que se espera obter/demonstrar ao final do
% trabalho (TCC1 + TCC2); não é necessário ter resultados definitivos nesta
% etapa.

Como resultado deste trabalho, espera-se identificar e comparar as principais abordagens utilizadas para identificação e estabelecimento da identidade de agentes de IA. A análise deverá considerar os requisitos levantados, potenciais aplicações e limitações das abordagens encontradas. Deve ser possível comparar entre as abordagens os seguintes aspectos:
\begin{itemize}
	\item diferenciar a identidade do usuário da identidade do agente que atua em seu nome;
	\item garantir a rastreabilidade das ações realizadas;
	\item permitir a operação entre diferentes domínios;
	\item gerenciar o ciclo de vida dos agentes.
\end{itemize}

Diferenciar a identidade do usuário da identidade do agente é fundamental para permitir que o agente seja reconhecido como uma entidade distinta, mesmo quando atua em nome do usuário, garantindo que o serviço que recebe a solicitação possa identificar a origem da ação, sem entregar todos os privilégios do usuário para o agente.

A rastreabilidade deve garantir a auditoria das ações realizadas pelo agente, sabendo exatamente qual ação foi realizada por qual agente, em qual contexto e em nome de qual usuário, permitindo a identificação de falhas, comportamentos inesperados. É esperado a rastreabilidade de múltiplos agentes atuando em nome de um mesmo usuário, assim como uma cadeia de delegação para sistemas com delegação entre agentes.

A operação entre diferentes domínios deve permitir que o agente se comunique em diferentes sistemas, serviços e organizações, ou com outros agentes. Onde é necessário que verificação da identidade e reconhecimento de confiança com o usuário que delegou a ação.

Gerenciamento do ciclo de vida dos agentes deve permitir que o agente seja criado, e atue apenas dentro do tempo necessário para cumprir a tarefa delegada, sendo encerrado após a conclusão da tarefa, serão estudados modelos persistentes, onde o agente pode ser reutilizado para outras tarefas, tendo um ciclo de vida mais longo, e modelos efêmeros, onde o agente é criado para uma Unica tarefa e é encerrado após a conclusão, tendo sua identidade e permissões descartadas.

\section{Considerações}
% Máximo de 200 palavras. Síntese da relevância do tema e do que se espera
% contribuir com o trabalho.
% Só aparecem aqui as referências efetivamente citadas no texto, com
% \cite{chave} ou \textcite{chave}.

A crescente adoção de agentes autônomos em diversos setores, desde assistentes pessoais até tomadas de decisões em sistemas críticos, torna essencial a implementação de mecanismos de identidade e acesso, para garantir a minimização de falhas e garantir a auditoria caso haja algum problema. 

O trabalho busca contribuir para a discussão sobre gestão de identidade aplicada à IA agêntica por meio da sistematização dos requisitos relacionados à autenticação, autorização, auditoria, bem como análise de modelos existentes para adequação desses modelos aos requisitos identificados. Dessa forma, buscar um identificador adequado, ver como sera construida a identidade digital do agente, o ciclo de vida necessário para a ação realizada e a resolução do identificador do agente em diferentes domínios de confiança. O trabalho pretende contribuir para a construção de um ecossistema agêntico mais seguro e confiável, onde a interação entre agentes e usuários seja transparente, rastreável e controlada.


\printbibliography[heading=bibintoc,title={Referências}]

\end{document}
